% !PRINTED 283
\weberpage{283}

%<<CONT
form. For it has already been shown earlier that the group $L_7$ contains divisors of degree 21 whose index is $= 8$, and which are of particular importance precisely for our purpose. Among the conjugate divisors of index 8, however, the simplest and the most convenient for calculation is the one consisting of all substitutions

\weq{5}{\begin{pmatrix} a, & b \\ 0, & a^{-1} \end{pmatrix} \quad (\mathrm{mod}\ 7)}

and we wish to arrange the transformation so that $\chi$ appears in this form. This we achieve by transformation of the substitutions (4) by means of $\begin{pmatrix} 1, & 1 \\ -1, & 0 \end{pmatrix}$, and thereby there results from (4)

\weq{6}{\chi = \begin{pmatrix} 2, & 3 \\ 0, & -3 \end{pmatrix}, \ \omega = \begin{pmatrix} 1, & -3 \\ 3, & -1 \end{pmatrix}, \ \Theta = \begin{pmatrix} 1, & 3 \\ -2, & 2 \end{pmatrix}.}

We now further assemble, for a better overview, the whole octahedral group into a table, whose six rows contain the elements $\Theta^\lambda, \omega\Theta^\lambda, \chi\Theta^\lambda, \chi\omega\Theta^\lambda, \chi^2\Theta^\lambda, \chi^2\omega\Theta^\lambda$ for $\lambda = 0, 1, 2, 3$:

\weq{7}{\begin{gathered}
\begin{pmatrix} 1, & 0 \\ 0, & 1 \end{pmatrix}, \ \begin{pmatrix} 1, & 3 \\ -2, & 2 \end{pmatrix}, \ \begin{pmatrix} 2, & 2 \\ 1, & -2 \end{pmatrix}, \ \begin{pmatrix} 2, & -3 \\ 2, & 1 \end{pmatrix};\\
\begin{pmatrix} 1, & -3 \\ 3, & -1 \end{pmatrix}, \ \begin{pmatrix} 0, & 3 \\ 2, & 0 \end{pmatrix}, \ \begin{pmatrix} -1, & 1 \\ -2, & 1 \end{pmatrix}, \ \begin{pmatrix} 3, & 1 \\ -3, & -3 \end{pmatrix};\\
\begin{pmatrix} 2, & 3 \\ 0, & -3 \end{pmatrix}, \ \begin{pmatrix} 3, & -2 \\ -1, & 1 \end{pmatrix}, \ \begin{pmatrix} 0, & 2 \\ 3, & 1 \end{pmatrix}, \ \begin{pmatrix} 3, & -3 \\ 1, & -3 \end{pmatrix};\\
\begin{pmatrix} 3, & 2 \\ 2, & -3 \end{pmatrix}, \ \begin{pmatrix} 1, & 1 \\ -1, & 0 \end{pmatrix}, \ \begin{pmatrix} 1, & 2 \\ 1, & 3 \end{pmatrix}, \ \begin{pmatrix} -3, & 0 \\ 2, & 2 \end{pmatrix};\\
\begin{pmatrix} 3, & 3 \\ 0, & -2 \end{pmatrix}, \ \begin{pmatrix} -3, & 1 \\ -3, & 3 \end{pmatrix}, \ \begin{pmatrix} 2, & 0 \\ -2, & -3 \end{pmatrix}, \ \begin{pmatrix} -2, & 1 \\ 3, & -2 \end{pmatrix};\\
\begin{pmatrix} -2, & 2 \\ 1, & 2 \end{pmatrix}, \ \begin{pmatrix} -1, & 2 \\ 3, & 0 \end{pmatrix}, \ \begin{pmatrix} 2, & 1 \\ 3, & 2 \end{pmatrix}, \ \begin{pmatrix} 0, & 1 \\ -1, & -1 \end{pmatrix}.
\end{gathered}}
To obtain the whole group $L_7$ we must still add an element of degree $7$, and for this we choose

\weq{8}{\tau = \begin{pmatrix} 1, & 1 \\ 0, & 1 \end{pmatrix}.}
